\chapter{Funções trigonométricas}\label{ch:g12:trigo}

As funções $\cos$ e $\sin$ transformam a geometria do círculo em análise.
Definidas enrolando a reta real em torno do círculo trigonométrico, elas
são o arquétipo das funções periódicas, e suas \omterm{def:g12:deriv:derivative}{derivadas} fazem delas
soluções da \omterm{def:g10:algebra:equation}{equação} das oscilações $y'' = -y$.

\section{O círculo trigonométrico}

\begin{definition}[Cosseno e seno]\label{def:g12:trigo:cossin}
Seja $\mathcal{C}$ o círculo de raio 1 centrado na origem de um \omterm{def:g10:coordgeom:system}{sistema de
coordenadas} \omterm{def:g10:coordgeom:system}{ortonormal}. A cada real $t$ associe o ponto $M(t)$ obtido
enrolando um comprimento $\abs{t}$ ao longo de $\mathcal{C}$ a partir do
ponto $I(1,0)$, no sentido anti-horário se $t \geq 0$ e no horário caso
contrário. Então \emph{$\cos t$ e $\sin t$ são as coordenadas de
$M(t)$}\index{cosseno}\index{seno}:
\[
M(t) = (\cos t, \sin t).
\]
\end{definition}

\begin{omfigure}
\begin{tikzpicture}[scale=2.2]
  \draw[->] (-1.25,0) -- (1.35,0) node[right] {$x$};
  \draw[->] (0,-1.25) -- (0,1.35) node[above] {$y$};
  \draw (0,0) circle (1);
  \coordinate (O) at (0,0);
  \coordinate (M) at (50:1);
  \draw[thick,omDef] (O) -- (M);
  \fill (M) circle (0.6pt) node[above right] {$M(t)$};
  \fill (1,0) circle (0.6pt) node[below right] {$I$};
  \draw[dashed] (M) -- (M |- O) node[below,pos=1.05] {$\cos t$};
  \draw[dashed] (M) -- (M -| O) node[left,pos=1.02] {$\sin t$};
  \draw[->,omThm,thick] (0.35,0) arc (0:50:0.35);
  \node[omThm] at (25:0.5) {$t$};
\end{tikzpicture}
\end{omfigure}

Como o círculo tem comprimento $2\pi$, o ponto $M(t)$ não muda quando $t$
aumenta de $2\pi$; e, como $M(t)$ está no círculo de raio $1$:

\begin{proposition}[Identidades fundamentais]\label{prop:g12:trigo:identities}
Para todos $t \in \R$ e $k \in \Z$:
\[
\cos^2 t + \sin^2 t = 1, \qquad
\cos(t + 2k\pi) = \cos t, \qquad
\sin(t + 2k\pi) = \sin t,
\]
\[
\cos(-t) = \cos t, \qquad \sin(-t) = -\sin t,
\]
\[
\cos(\pi - t) = -\cos t, \quad \sin(\pi - t) = \sin t, \quad
\cos\left(\tfrac{\pi}{2} - t\right) = \sin t, \quad
\sin\left(\tfrac{\pi}{2} - t\right) = \cos t .
\]
\end{proposition}

\begin{proof}
A primeira identidade é a \omterm{def:g10:algebra:equation}{equação} do círculo de raio $1$; as demais
exprimem as simetrias do círculo: $M(-t)$ é o reflexo de $M(t)$ no eixo
$x$, $M(\pi - t)$ no eixo $y$, e $M(\frac\pi2 - t)$ na reta $y = x$.
\end{proof}

Os valores que se devem saber de cor:
\begin{center}
\begin{tabular}{c|ccccc}
$t$ & $0$ & $\dfrac{\pi}{6}$ & $\dfrac{\pi}{4}$ & $\dfrac{\pi}{3}$ & $\dfrac{\pi}{2}$\\[6pt]
\hline\rule{0pt}{14pt}%
$\cos t$ & $1$ & $\dfrac{\sqrt3}{2}$ & $\dfrac{\sqrt2}{2}$ & $\dfrac12$ & $0$\\[6pt]
$\sin t$ & $0$ & $\dfrac12$ & $\dfrac{\sqrt2}{2}$ & $\dfrac{\sqrt3}{2}$ & $1$
\end{tabular}
\end{center}

\begin{proposition}[Fórmulas de adição]\label{prop:g12:trigo:addition}
Para todos $a, b \in \R$:
\begin{align*}
\cos(a+b) &= \cos a \cos b - \sin a \sin b, &
\sin(a+b) &= \sin a \cos b + \cos a \sin b,\\
\cos(a-b) &= \cos a \cos b + \sin a \sin b, &
\sin(a-b) &= \sin a \cos b - \cos a \sin b.
\end{align*}
Em particular, $\cos 2a = 2\cos^2 a - 1 = 1 - 2\sin^2 a$ e
$\sin 2a = 2 \sin a \cos a$.
\end{proposition}

\begin{proof}
Os \omterm{def:g11:vect:vector}{vetores} $\vect{u} = (\cos a, \sin a)$ e $\vect{v} = (\cos b, \sin b)$
são unitários e formam um ângulo $a - b$; seu \omterm{def:g11:scal:dot}{produto escalar} vale tanto
$\cos a\cos b + \sin a \sin b$ (\omterm{def:g10:coordgeom:system}{coordenadas}) quanto
$\norm{\vect u}\,\norm{\vect v}\cos(a-b) = \cos(a - b)$ (geometria), o que
é a terceira fórmula. Trocar $b$ por $-b$ dá a primeira; as fórmulas do
\omterm{def:g12:trigo:cossin}{seno} seguem usando $\sin x = \cos(\frac\pi2 - x)$. As fórmulas do arco
duplo são o caso $b = a$ combinado com $\cos^2 + \sin^2 = 1$.
\end{proof}

\section{Estudo analítico}

\begin{lemma}[Limite fundamental]\label{lem:g12:trigo:sinxx}
\[
\lim_{t \to 0} \frac{\sin t}{t} = 1
\qquad\text{e}\qquad
\lim_{t \to 0} \frac{\cos t - 1}{t} = 0 .
\]
\end{lemma}

\begin{proof}
Para $0 < t < \frac{\pi}{2}$, compare três áreas no círculo de raio $1$: o
triângulo $OIM(t)$, o setor circular $OIM(t)$ e o triângulo retângulo de
base $OI$ e altura $\tan t$:
\[
\frac{\sin t}{2} \leq \frac{t}{2} \leq \frac{\tan t}{2}.
\]

\begin{omfigure}
\begin{tikzpicture}[scale=2.2]
  \draw[->] (-0.25,0) -- (1.5,0) node[right] {$x$};
  \draw[->] (0,-0.25) -- (0,1.45) node[above] {$y$};
  \fill[omDef!20] (0,0) -- (1,0) -- (50:1) -- cycle;
  \draw[gray, thin] (-10:1) arc (-10:80:1);
  \draw[omProp, thick] (0,0) -- (1,0) -- (1,1.1918) -- cycle;
  \draw[omDef] (0,0) -- (1,0) -- (50:1) -- cycle;
  \draw[omThm, thick] (1,0) arc (0:50:1);
  \draw[dashed] (50:1) -- (0.643,0);
  \node[left, font=\small] at (0.68,0.35) {$\sin t$};
  \node[omProp, right, font=\small] at (1,0.62) {$\tan t$};
  \node[omThm, font=\small] at (20:0.5) {$t$};
  \fill (50:1) circle (0.6pt);
  \node at (58:1.22) {$M(t)$};
  \fill (1,0) circle (0.6pt) node[below right] {$I$};
  \node[below left] at (0,0) {$O$};
\end{tikzpicture}

{\small As três áreas encaixadas: o triângulo $OIM$ (azul, área
$\frac{\sin t}{2}$), o setor circular (delimitado pelo arco vermelho,
área $\frac{t}{2}$) e o triângulo retângulo de altura $\tan t$ (laranja,
área $\frac{\tan t}{2}$).}
\end{omfigure}

A primeira desigualdade dá $\frac{\sin t}{t} \leq 1$; a segunda dá
$\cos t \leq \frac{\sin t}{t}$. Quando $t \to 0^+$, $\cos t \to 1$
(\omterm{def:g12:limcont:continuity}{continuidade}), e o teorema do confronto dá $\frac{\sin t}{t} \to 1$; a
paridade estende isso a $t \to 0$. Para o segundo limite,
\[
\frac{\cos t - 1}{t} = \frac{\cos^2 t - 1}{t(\cos t + 1)}
= -\frac{\sin t}{t}\cdot\frac{\sin t}{\cos t + 1}
\xrightarrow[t\to0]{} -1 \cdot \frac{0}{2} = 0. \qedhere
\]
\end{proof}

\begin{theorem}[Derivadas do seno e do cosseno]\label{thm:g12:trigo:derivatives}
As funções $\sin$ e $\cos$ são deriváveis em $\R$, com
\[
(\sin)' = \cos, \qquad (\cos)' = -\sin .
\]
Mais geralmente, $(\sin(ax+b))' = a\cos(ax+b)$ e
$(\cos(ax+b))' = -a\sin(ax+b)$.
\end{theorem}

\begin{proof}
Pelas fórmulas de adição,
\[
\frac{\sin(x+h) - \sin x}{h}
= \sin x\,\frac{\cos h - 1}{h} + \cos x\,\frac{\sin h}{h}
\xrightarrow[h\to0]{} \sin x \cdot 0 + \cos x \cdot 1 = \cos x,
\]
usando o \cref{lem:g12:trigo:sinxx}. O cálculo para $\cos$ é idêntico, e a
forma geral segue da regra da cadeia.
\end{proof}

\begin{proposition}[Variação]\label{prop:g12:trigo:variations}
Em $\intcc{0}{\pi}$, $\cos$ decresce de $1$ a $-1$. Em
$\intcc{-\frac\pi2}{\frac\pi2}$, $\sin$ cresce de $-1$ a $1$. As duas
funções são $2\pi$-periódicas e têm \omterm{def:g10:functions:function}{imagem} $\intcc{-1}{1}$; $\cos$ é par e
$\sin$ é ímpar.
\end{proposition}

\begin{proof}
Em $\intoo{0}{\pi}$, $(\cos)' = -\sin < 0$, pois o ponto $M(t)$ tem ali
ordenada positiva; analogamente, $(\sin)' = \cos > 0$ em
$\intoo{-\frac\pi2}{\frac\pi2}$. A periodicidade e a paridade vêm da
\cref{prop:g12:trigo:identities}.
\end{proof}

\begin{omfigure}
\begin{tikzpicture}
\begin{axis}[
  width=12cm, height=4.5cm,
  axis lines=middle,
  xmin=-7, xmax=7, ymin=-1.4, ymax=1.4,
  xtick={-6.2832,-3.1416,3.1416,6.2832},
  xticklabels={$-2\pi$,$-\pi$,$\pi$,$2\pi$},
  ytick={-1,1},
  samples=200, domain=-7:7]
  \addplot[omDef,thick] {cos(deg(x))};
  \addplot[omThm,thick] {sin(deg(x))};
\end{axis}
\end{tikzpicture}

{\small As curvas de $\cos$ (azul) e de $\sin$ (vermelho).}
\end{omfigure}

\begin{method}[Resolver $\cos x = a$ e $\sin x = a$]\label{met:g12:trigo:equations}
Para $a \in \intcc{-1}{1}$, encontre uma solução particular $\alpha$ (pela
tabela de valores ou por uma calculadora). Então todas as soluções são:
\[
\cos x = \cos\alpha \iff x = \alpha + 2k\pi \text{ ou } x = -\alpha + 2k\pi,
\quad k \in \Z;
\]
\[
\sin x = \sin\alpha \iff x = \alpha + 2k\pi \text{ ou } x = \pi - \alpha + 2k\pi,
\quad k \in \Z.
\]
Para inequações, localize os arcos solução no círculo trigonométrico e leia
os \omterm{def:g10:numbers:interval}{intervalos} dentro de um período.
\end{method}

\begin{example}\label{ex:g12:trigo:equation}
Resolva $\cos x = \frac12$ em $\intoc{-\pi}{\pi}$: uma solução particular é
$\frac\pi3$, de modo que as soluções são $x = \frac\pi3$ e
$x = -\frac\pi3$. Em todo o $\R$: $x = \pm\frac\pi3 + 2k\pi$, $k \in \Z$.
\end{example}

\section{Exercícios}

\begin{exercise}[$\star$]\label{exo:g12:trigo:1}
Calcule $\cos\frac{2\pi}{3}$, $\sin\frac{5\pi}{6}$, $\cos\frac{7\pi}{4}$ e
$\sin\left(-\frac{\pi}{3}\right)$ usando as identidades da
\cref{prop:g12:trigo:identities}.
\end{exercise}

\begin{exercise}[$\star$]\label{exo:g12:trigo:2}
Derive $f(x) = \sin^2 x$, $g(x) = \cos(3x) + x\sin x$ e
$h(x) = \dfrac{\sin x}{2 + \cos x}$.
\end{exercise}

\begin{exercise}[$\star$]\label{exo:g12:trigo:3}
Resolva em $\intoc{-\pi}{\pi}$ e depois em $\R$:
\[
\text{(a) } \sin x = \frac{\sqrt3}{2};
\qquad
\text{(b) } \cos\left(2x\right) = \frac{\sqrt2}{2} .
\]
\end{exercise}

\begin{exercise}[$\star\star$]\label{exo:g12:trigo:4}
Usando as fórmulas de adição, calcule o valor exato de
$\cos\frac{\pi}{12}$ e de $\sin\frac{\pi}{12}$. (Sugestão:
$\frac{\pi}{12} = \frac{\pi}{3} - \frac{\pi}{4}$.)
\end{exercise}

\begin{exercise}[$\star\star$]\label{exo:g12:trigo:5}
Resolva em $\intcc{0}{2\pi}$ a inequação $2\sin^2 x - \sin x - 1 \geq 0$.
(Sugestão: fatore o trinômio em $\sin x$.)
\end{exercise}

\begin{exercise}[$\star\star$]\label{exo:g12:trigo:6}
Calcule os limites
\[
\lim_{x\to0} \frac{\sin 3x}{x}, \qquad
\lim_{x\to0} \frac{1 - \cos x}{x^2}, \qquad
\lim_{x\to+\infty} x \sin\frac{1}{x}.
\]
\end{exercise}

\begin{exercise}[$\star\star$]\label{exo:g12:trigo:7}
Estude a \omterm{def:g11:func:function}{função} $f(x) = x + \cos x$ em $\intcc{0}{2\pi}$: variação e
\omterm{def:g12:deriv:extremum}{extremos}. Mostre que $f$ é \omterm{def:g12:seq:monotonic}{crescente} em $\R$ embora $f'$ se anule em
infinitos pontos.
\end{exercise}

\begin{exercise}[$\star\star\star$]\label{exo:g12:trigo:8}
Seja $f(x) = A\cos x + B \sin x$ com $(A,B) \neq (0,0)$.
\begin{enumerate}
  \item Mostre que $f$ pode ser escrita como $f(x) = R\cos(x - \varphi)$
        com $R = \sqrt{A^2 + B^2}$ e um $\varphi$ adequado.
  \item Deduza o \omterm{def:g10:functions:extrema}{máximo} e o \omterm{def:g10:functions:extrema}{mínimo} de $f$ e resolva
        $\cos x + \sin x = 1$ em $\intoc{-\pi}{\pi}$.
\end{enumerate}
\end{exercise}

\begin{exercise}[$\star\star\star$]\label{exo:g12:trigo:9}
Mostre que $\sin$ e $\cos$ satisfazem ambas a \omterm{def:g10:algebra:equation}{equação} diferencial
$y'' = -y$. Reciprocamente, seja $y$ uma solução de $y'' = -y$ com
$y(0) = a$ e $y'(0) = b$; mostre que $y = a\cos + b\sin$. (Sugestão:
considere $g = (y - a\cos - b\sin)$ e a ``energia''
$E = g^2 + (g')^2$.)
\end{exercise}

\section{Problema: O segredo do pêndulo}

\begin{problem}[{Problema de fim de semana --- um limite geométrico
move as derivadas do seno e do cosseno e, a partir delas, toda a
oscilação: molas, pêndulos, batimentos e o metro que quase foi}]
\label{pb:g12:trigo:1}
Todo relógio que já bateu, toda corda que já soou obedece à mesma
matemática: uma força restauradora proporcional ao deslocamento,
portanto a \omterm{def:g10:algebra:equation}{equação} $y'' = -\omega^2 y$, portanto \omterm{def:g12:trigo:cossin}{cossenos}. Na base
de tudo está um limite de aparência inocente,
$\frac{\sin x}{x} \to 1$ --- prometido no ano 9, usado pelo
\cref{thm:g12:trigo:derivatives} e demonstrado aqui com um
desenho. Este problema sobe desse limite até a oscilação de um
pêndulo de um metro e explica por que o metro, por um fio, não é
definido por ele.

\medskip
\noindent\textbf{Parte I --- O limite fundamental.}
\begin{enumerate}
  \item Aqueça numericamente: calcule $\frac{\sin x}{x}$ para
        $x = 0.1$ e $x = 0.01$ (em \omterm{def:g11:trigo:radian}{radianos}!).
  \item O sanduíche, no círculo de raio $1$ com
        $0 < x < \frac\pi2$: exprima por geometria elementar as
        áreas (i) do triângulo de vértices $O$, $A(1,0)$ e o
        ponto $M(\cos x, \sin x)$ do círculo; (ii) do setor
        circular $OAM$; (iii) do triângulo retângulo $OAT$, com
        $T(1, \tan x)$.
  \item Do sanduíche de áreas deduza $\sin x < x < \tan x$ e
        depois $\cos x < \frac{\sin x}{x} < 1$, e conclua pelo
        teorema do confronto que
        $\lim_{x \to 0} \frac{\sin x}{x} = 1$ (a paridade trata
        de $x < 0$).
  \item Deduza $\lim_{x \to 0} \frac{1 - \cos x}{x^2} =
        \frac12$ (multiplique pelo conjugado).
  \item Explique por que esse limite é o motor do
        \cref{thm:g12:trigo:derivatives} (quanto vale
        $\sin'(0)$?) e por que ele também salda uma dívida
        antiga: em que sentido os \emph{\omterm{def:g11:trigo:radian}{radianos}} são a única
        unidade de ângulo para a qual $\sin' = \cos$ vale sem
        constante?
\end{enumerate}

\medskip
\noindent\textbf{Parte II --- Movimento harmônico.}
\begin{enumerate}[resume]
  \item Verifique pela regra da cadeia que
        $y(t) = R\cos(\omega t - \varphi)$ satisfaz
        $y'' = -\omega^2 y$.
  \item Uma massa presa a uma mola obedece a $y'' = -4y$, com
        $y(0) = 3$ e $y'(0) = 8$. Admitindo (do
        \cref{exo:g12:trigo:9}) que todas as soluções são
        $A\cos 2t + B\sin 2t$, determine $A$ e $B$.
  \item Reescreva a solução como $R\cos(2t - \varphi)$
        (\cref{exo:g12:trigo:8}): dê a amplitude $R$, o período e
        a fase $\varphi$ (três casas decimais).
  \item Em que instante a massa passa pela primeira vez pela
        posição de equilíbrio $y = 0$?
  \item O pêndulo: uma massa presa a um fio de comprimento $L$
        obedece \emph{exatamente} a
        $\theta'' = -\frac gL \sin\theta$ --- insolúvel por
        funções elementares. Para oscilações pequenas, substitua
        $\sin\theta$ por $\theta$ (a Parte I justifica!) e deduza
        o período $T = 2\pi\sqrt{\frac Lg}$. Calcule $T$ para
        $L = 1$ m ($g = 9.81$): a famosa descoberta de Galileu
        --- de que a fórmula \emph{não} depende?
  \item O pêndulo de segundos: que comprimento dá $T = 2$ s
        exatamente? A Academia francesa de 1791 hesitou entre
        esse comprimento e a décima milionésima parte do quarto
        de meridiano para definir o \emph{metro} --- por quantos
        milímetros os dois candidatos diferem?
\end{enumerate}

\medskip
\noindent\textbf{Parte III --- As fórmulas de adição em serviço.}
\begin{enumerate}[resume]
  \item Mostre, com as fórmulas de adição
        (\cref{prop:g12:trigo:addition}), que
        $\cos\left(x + \frac\pi2\right) = -\sin x$ e
        $\sin\left(x + \frac\pi2\right) = \cos x$: derivar no
        círculo é dar um quarto de volta --- confira a coerência
        com o \cref{thm:g12:trigo:derivatives}.
  \item Deduza a identidade produto-em-soma
        $\cos a \cos b = \frac12\left(\cos(a - b) +
        \cos(a + b)\right)$ e a linearização
        $\cos^2 x = \frac{1 + \cos 2x}{2}$ (o capítulo de
        integração lhe agradecerá).
  \item Batimentos: deduza
        $\sin p + \sin q = 2 \sin\frac{p + q}{2}
        \cos\frac{p - q}{2}$ e aplique isso a dois diapasões de
        $440$ e $444$ Hz: que som o ouvido escuta e quantas
        pulsações por segundo? (É assim que as orquestras
        afinam.)
  \item Deduza a identidade do arco triplo
        $\cos 3x = 4\cos^3 x - 3\cos x$ (escreva $3x = 2x + x$).
        Fazendo $x = 20^\circ$, que \omterm{def:g10:algebra:equation}{equação} cúbica
        $\cos 20^\circ$ tem de satisfazer? (Que essa cúbica não
        possa ser resolvida só com raízes quadradas é
        precisamente a razão pela qual a trissecção do ângulo
        derrotou régua e compasso por dois milênios --- a
        história completa é contada nos volumes de graduação.)
  \item Uma corrente alternada vale
        $I(t) = 3\cos(100\pi t) + 3\sqrt3 \sin(100\pi t)$.
        Reescreva-a como $R\cos(100\pi t - \varphi)$: corrente de
        pico e fase?
\end{enumerate}

\medskip
\noindent\textbf{Parte IV --- O mundo que oscila.}
\begin{enumerate}[resume]
  \item Períodos: qual é o (menor) período de $\sin^2 t$ (use a
        questão 13)? E de $\sin(2t) + \sin(3t)$?
  \item Um pêndulo real amortece:
        $y(t) = \eu^{-t/10}\cos(2\pi t)$. Descreva o movimento,
        calcule o valor da envoltória em $t = 10$ (aparece um
        velho conhecido, \cref{pb:g12:exp:1}) e nomeie o capítulo
        em que se resolvem \omterm{def:g10:algebra:equation}{equações} com soluções desse tipo.
  \item Por que os \omterm{def:g12:trigo:cossin}{senos} e \omterm{def:g12:trigo:cossin}{cossenos} reinam sobre toda vibração?
        Dê o dicionário física-matemática em duas frases (força
        restauradora proporcional ao deslocamento $\rightarrow$
        que \omterm{def:g10:algebra:equation}{equação} $\rightarrow$ que soluções) e diga o que o
        \cref{exo:g12:trigo:9} acrescenta a essa afirmação.
  \item Final: o arco do capítulo em quatro passos --- um limite
        geométrico, duas \omterm{def:g12:deriv:derivative}{derivadas}, uma \omterm{def:g10:algebra:equation}{equação} diferencial, toda
        a oscilação. E o asterisco honesto: para oscilações
        grandes, o período verdadeiro do pêndulo envolve uma
        integral elíptica, calculada à velocidade da luz
        pela\,\dots \omterm{def:g11:stat:mean}{média} aritmético-geométrica de Gauss, a
        sobremesa do \cref{pb:g12:seq:1}. Os fios da série
        amarram-se sozinhos.
\end{enumerate}
\end{problem}
